\documentclass[11pt]{amsart}
\usepackage[utf8]{inputenc}
\usepackage[margin=1in, headheight=14.5pt]{geometry}
\usepackage{math-hw}
\pagestyle{fancy}
\fancyhead[L]{Math 918}
%\fancyhead[C]{MATH 665}
\fancyhead[R]{HW 3}

\usepackage[margin=1in, headheight=14.5pt]{geometry}
\usepackage{parskip}

%%%%%%%%%%%%%%%%%%%%%%%%%%%%%%%%%%
% fonts
\usepackage{calligra}
\usepackage{mathrsfs}
\usepackage[mathscr]{euscript}
\usepackage{bbm}

\usepackage{enumitem}
\usepackage{hyperref}
%\usepackage[colorlinks,  citecolor=green, linkcolor = black, urlcolor = cyan]{hyperref} % hyperlinks for cross-referencing; colorlinks is recommended because otherwise the boxes obscure the roman numerals; linkcolor is changed to blue after the table of contents (see firstpages.tex). As of 2021, red is not allowed by Rackham.
\usepackage[nameinlink]{cleveref}
%\usepackage{backref}
\usepackage{url} % allow hyperlink in reference
\usepackage[style=alphabetic,url=false,backref=true]{biblatex}
\addbibresource{zotero-export.bib}
\setlength{\droptitle}{-4em}


\newtoggle{solution}
\toggletrue{solution}
%\togglefalse{solution}


\iftoggle{solution}{
\includeversion{solnenv}
\includeversion{commenv}
}{
\excludeversion{solnenv}
\excludeversion{commenv}
}





\title{Homework 3}
%\author{Anna Brosowsky}
\date{\vspace{-4em}}

\begin{document}


\maketitle\thispagestyle{fancy}

Remember you are allowed to discuss with classmates (or an AI tool), but that you need to tell me who/what you discussed with + the final submitted writeup should be your own work.


\begin{enumerate}[itemsep=7em]
\item Suppose that $S=R_1\times R_2$, and that $R_1,R_2$, and $S$ are all $F$-finite\footnote{In fact this result holds without the $F$-finiteness hypothesis (and both $R_i$ F-finite automatically implies $S$ F-finite), but I am including it since for us $F$-finiteness is part of the definition of strong $F$-regularity}. 
Prove that $S$ is strongly $F$-regular if and only if both $R_1$ and $R_2$ are strongly $F$-regular\footnote{Hint: There are multiple ways to approach this. One approach is mentioned in the proof of Theorem~3.9 of [MP], which you can look at if you get stuck!}.



\item (Exercise 13 of [MP]) Let $R\to S$ be a faithfully flat extension of $F$-finite rings. 
\begin{enumerate}[itemsep=1em]
\item Let $M,N$ be finitely generated $R$-modules, and let $\psi:M\to N$ be an $R$-module map. Prove that $\psi$ is surjective if and only if $\psi\otimes_R \id_S: M\otimes_R S \to N\otimes_R S$ is surjective.


\item Prove that if $c$ is a non-zerodivisor on $R$, then $c$ is also a non-zerodivisor on $S$.

\item Prove that if $S$ is strongly $F$-regular, then so is $R$.

\end{enumerate}


\item (Glassbrenner's Criterion). Let $(S, \mf m, k)$ be an F -finite regular local ring of prime
characteristic $p > 0$ and let $I \subseteq S$ be an ideal\footnote{Note: Just like Fedder, Glassbrenner's criterion also holds if $S$ is a standard graded ring over an $F$-finite field,$\mf m$ is the homogeneous maximal ideal, and $I$ is a homogenous ideal}. Let $R = S/I$. 
\begin{enumerate}[itemsep=1em]
\item Let $c\in S$. Prove that $R$ is $e$-F-split along $c$ if and only if $c(I^{[p^e]}:I)\not\subseteq \mf m^{[p^e]}$.\footnote{Hint: Use results from Jan 29th \& Feb 3rd! And try to modify the proof of Fedder's criterion.}

\item Conclude that the following are equivalent: 
    \begin{enumerate}
    \item $R$ is strongly $F$-regular.
    \item For every $c \in S$ not in any minimal prime of $I$, there exists $e > 0$ such that $c\not\in\mf m^{[p^e]}:(I^{[p^e]} :I)$.
    \item $\bigcap_{e> 0}\mf m^{[p^e]}:(I^{[p^e]}:I)$ is contained in the union of the minimal primes.\footnote{Fun fact: By prime avoidance, you could even restate this a fourth way and require that the intersection by contained in a single minimal prime!}
    \end{enumerate}
\end{enumerate}



\newpage

\item Prove the following number theory results. Since (a) is a special case of (b), you only need to write up a solution to (b). However, you may find it helpful to think about the simpler part (a) first!
\begin{enumerate}[itemsep=1em]
\item (Lucas's Theorem) Write $n=\sum_{i=0}^d n_i p^i$ and $m = \sum_{i=0}^d m_i p^i$ where $0\le n_i,m_i<p$ for all $i$.\footnote{Hint: Both parts (a) \& (b) are still true if you allow the top coefficient to be even larger, i.e., if we relax the last inequality to just be $n_d,m_d,m_{d,i}\ge 0$. You may actually find it easier to prove this more general version!} Then
\[
\binom{n}{m} \equiv \prod_{i=0}^d \binom{n_i}{m_i} \mod p.
\]
Here we follow the convention that $\binom{a}{b} = 0$ if $b>a$.


\item (Multinomial Lucas's Theorem) Suppose that $n = m_1+\cdots + m_t$. Write $n=\sum_{i=0}^d n_ip^i$ and $m_j = \sum_{i=0}^d m_{i,j} p^i$ where $0\le n_i,m_{i,j}<p$ for all $i$.
Then
\[
\binom{n}{m_1,\ldots, m_t} \equiv \prod_{i=0}^d \binom{n_i}{m_{i,1},\ldots, m_{i,t}} \mod p.
\]
Here we follow the convention that $\binom{a}{b_1,\ldots, b_t}=0$ if $b_1+\cdots + b_t\neq a$.

\end{enumerate}


\item (Exercise 17 of [MP]) Let $k$ be an algebraically closed\footnote{This problem is also true with the weaker hypothesis that $k$ is $F$-finite. You would prove that from the algebraically closed case by using Problem 2 and the fact that if $R\to S$ is faithfully flat, so is $R/I\to S/IS$.} 
field and let 
\[R = k[ x_1, \ldots , x_d]/\langle x^n_1 +\cdots + x^n_d \rangle.\] 
Using Glassbrenner's criterion, prove each of the following facts:
\begin{enumerate}[itemsep=1em]
\item Show that $R$
is not strongly F -regular if $n \ge d \ge 2$.
\item Show that $R$ is strongly $F$-regular if $n < d$ and $p \gg 0$.\footnote{Hint: Use ``test element'' criterion (Thm 3.11 of [MP]; done in-class on Tuesday 02/24); an algebraic geometry criterion for non-singularity;  and use some number theory to make divisibility things nicer.}
\end{enumerate}





\end{enumerate}


\end{document}
